\documentclass[11pt,a4paper]{article}

\usepackage{fontsetup}
\usepackage[english, russian, ukrainian]{babel}


\usepackage{amsmath}
\usepackage{geometry}
\geometry{margin=2.5cm}
\usepackage{enumitem}
\usepackage{tcolorbox}
\usepackage{hyperref}

\title{Зауваження до конспекту лекцій С.~М.~Пономаренка\\«Електрика та магнетизм»}
\author{Аналіз стилю викладення}
\date{}

\begin{document}
\maketitle

\noindent
Наведено конкретні місця в конспекті Пономаренка, де автор пропускає проміжні викладки, вважаючи їх очевидними для читача. Це не критика~--- це особливість лекційного стилю, але для читача, який звик до покрокових підручників, ці місця можуть бути складними.

\section{Виведення поля точкового диполя (розділ 3.2)}

\textbf{Що написано:}
\[
\mathbf{E} = \frac{3(\mathbf{p} \cdot \mathbf{r})\mathbf{r}}{r^5} - \frac{\mathbf{p}}{r^3}
\]
і одразу~--- «напруженість поля точкового диполя зменшується з відстанню за законом \(E \sim 1/r^3\)».

\textbf{Що пропущено:}
\begin{itemize}[nosep]
    \item Проміжний крок розкладу \(1/|\mathbf{r} - \mathbf{r}_\pm|\) в ряд Тейлора за малим \(\ell/r\).
    \item Підстановка \(\mathbf{r}_\pm = \mathbf{r} \mp \boldsymbol{\ell}/2\) і збереження членів першого порядку.
    \item Алгебраїчне спрощення, де з'являється множник \(3\) і векторний добуток.
\end{itemize}

\textbf{Чому це проблема:} студент бачить готову формулу, але не розуміє, звідки взялася «3» і чому саме така комбінація векторів.

\section{Виведення потенціалу системи зарядів на далеких відстанях (розділ 3.4)}

\textbf{Що написано:}
\[
\frac{1}{|\mathbf{r} - \mathbf{r}_i|} \approx \frac{1}{r} + \frac{\mathbf{r} \cdot \mathbf{r}_i}{r^3} + \frac{3(\mathbf{r} \cdot \mathbf{r}_i)^2 - r^2 r_i^2}{2r^5}
\]

\textbf{Що пропущено:}
\begin{itemize}[nosep]
    \item Саме виведення цього розкладу. Автор просто подає його як «утримуючи в розкладі доданки до другого порядку малості».
    \item Проміжні кроки переходу від \(\varphi = \sum_i q_i/|\mathbf{r} - \mathbf{r}_i|\) до виразу з \(q\), \(\mathbf{p}\) і \(D_{\alpha\beta}\).
    \item Деталі введення тензора квадрупольного моменту \(D_{\alpha\beta} = \sum_i q_i(3r_{i\alpha}r_{i\beta} - r_i^2 \delta_{\alpha\beta})\).
\end{itemize}

\textbf{Чому це проблема:} для читача, який не звик до тензорних позначень, раптова поява \(D_{\alpha\beta}\) і символу Кронекера \(\delta_{\alpha\beta}\) виглядає як «магія».

\section{Виведення формули Лармора (розділ 15.5)}

\textbf{Що написано:}
\[
\frac{dW}{dt} = \frac{2q^2}{3c^3} \dot{\mathbf{v}}^2
\]
і одразу~--- «формулу називають формулою Лармора».

\textbf{Що пропущено:}
\begin{itemize}[nosep]
    \item Перехід від \(\ddot{\mathbf{p}} = -q\dot{\mathbf{v}}\) до \(dW/dt = 2\ddot{p}^2/(3c^3)\).
    \item Інтегрування вектора Пойнтінга по сфері~--- проміжний результат \(\int_0^\pi \sin^3\theta \, d\theta = 4/3\) подано, але без пояснення, чому саме такий інтеграл виникає.
    \item Усереднення за період для гармонічного осцилятора: \(\langle \cos^2 \omega t \rangle = 1/2\)~--- подано без виведення.
\end{itemize}

\textbf{Чому це проблема:} формула Лармора~--- одна з ключових у курсі, і для її розуміння важливо бачити, як саме виникає множник \(2/3\).

\section{Виведення хвильового рівняння для \(\mathbf{E}\) і \(\mathbf{H}\) (розділ 14.2)}

\textbf{Що написано:}
\[
\nabla^2 \mathbf{E} = \frac{1}{v^2} \frac{\partial^2 \mathbf{E}}{\partial t^2}, \quad v = \frac{c}{\sqrt{\varepsilon\mu}}
\]

\textbf{Що пропущено:}
\begin{itemize}[nosep]
    \item Деталі застосування операції \(\nabla \times\) до другого рівняння Максвелла.
    \item Використання тотожності \(\nabla \times \nabla \times \mathbf{E} = \nabla(\nabla \cdot \mathbf{E}) - \nabla^2 \mathbf{E}\)~--- подано, але без пояснення, чому \(\nabla \cdot \mathbf{E} = 0\).
    \item Перехід від \(\nabla \times \mathbf{H}\) до \(\partial \mathbf{E}/\partial t\) через четверте рівняння~--- один рядок, але без проміжних підстановок.
\end{itemize}

\textbf{Чому це проблема:} для читача, який не звик до векторних тотожностей, ці кроки неочевидні.

\section{Виведення теореми Пойнтінга (розділ 12.3)}

\textbf{Що написано:}
\[
-\frac{\partial w}{\partial t} = \nabla \cdot \boldsymbol{\Pi} + \mathbf{j} \cdot \mathbf{E}
\]
і далі~--- інтегральна форма.

\textbf{Що пропущено:}
\begin{itemize}[nosep]
    \item Проміжні кроки підстановки \(\mathbf{j} = \frac{c}{4\pi} \nabla \times \mathbf{H} - \frac{1}{c} \frac{\partial \mathbf{D}}{\partial t}\) у \(\mathbf{j} \cdot \mathbf{E}\).
    \item Використання тотожності \(\mathbf{E} \cdot \nabla \times \mathbf{H} = -\nabla \cdot [\mathbf{E} \times \mathbf{H}] + \mathbf{H} \cdot \nabla \times \mathbf{E}\)~--- подано, але без пояснення.
    \item Підстановка закону Фарадея \(\nabla \times \mathbf{E} = -\frac{1}{c} \frac{\partial \mathbf{B}}{\partial t}\) і матеріальних рівнянь~--- кілька рядків, але без коментарів.
\end{itemize}

\textbf{Чому це проблема:} теорема Пойнтінга~--- концептуально важлива, і її виведення варте детального опрацювання.

\section{Виведення телеграфних рівнянь (розділ 17.2)}

\textbf{Що написано:}
\[
-\frac{\partial u}{\partial z} = R_0 i + L_0 \frac{\partial i}{\partial t}, \quad -\frac{\partial i}{\partial z} = G_0 u + C_0 \frac{\partial u}{\partial t}
\]

\textbf{Що пропущено:}
\begin{itemize}[nosep]
    \item Деталі застосування другого правила Кірхгофа до елементарного відрізка \(dz\)~--- подано, але без пояснення, чому похідна \(\partial u/\partial z\) виникає саме так.
    \item Перехід від скінченного відрізка до границі \(dz \to 0\)~--- один рядок.
    \item Фізичний зміст кожного з погонних параметрів \(R_0, L_0, G_0, C_0\)~--- подано, але без прикладів.
\end{itemize}

\textbf{Чому це проблема:} для читача, який не звик до диференціальних рівнянь у частинних похідних, ці кроки неочевидні.

\section{Виведення формул Френеля (розділ 14.5)}

\textbf{Що написано:}
\[
E'_\perp = -\frac{\sin(i-r)}{\sin(i+r)} E_\perp, \quad E''_\perp = \frac{2\sin r \cos i}{\sin(i+r)} E_\perp
\]

\textbf{Що пропущено:}
\begin{itemize}[nosep]
    \item Проміжні кроки підстановки граничних умов для \(\mathbf{E}\) і \(\mathbf{H}\).
    \item Використання закону Снелла для спрощення виразів.
    \item Алгебраїчні перетворення, де з'являються синуси кутів \(i\) і \(r\).
\end{itemize}

\textbf{Чому це проблема:} формули Френеля~--- ключові для оптики, і їх виведення варте детального розгляду.

\section{Тензор електромагнітного поля (розділ 13.3)}

\textbf{Що написано:}
\[
F_{\mu\nu} = \partial_\mu A_\nu - \partial_\nu A_\mu
\]
і одразу~--- матриця \(4 \times 4\) з компонентами \(\mathbf{E}\) і \(\mathbf{B}\).

\textbf{Що пропущено:}
\begin{itemize}[nosep]
    \item Пояснення, чому саме така комбінація похідних дає тензор.
    \item Виведення компонент \(F_{0i} = E_i\), \(F_{12} = -B_z\) тощо~--- подано, але без проміжних кроків.
    \item Перехід від тензора до законів перетворення полів~--- подано, але без деталей.
\end{itemize}

\textbf{Чому це проблема:} тензорний запис~--- це «вища математика» для багатьох студентів, і без проміжних кроків він виглядає як абстракція.

\section{Загальна закономірність}

Пономаренко пропускає проміжні викладки в місцях, де:
\begin{enumerate}[nosep]
    \item \textbf{Використовуються векторні тотожності}~--- автор вважає, що читач їх знає.
    \item \textbf{Використовуються тензори й 4-вектори}~--- автор вважає, що читач уже знайомий з цим апаратом.
    \item \textbf{Виконуються алгебраїчні спрощення}~--- автор вважає їх рутинними.
    \item \textbf{Використовуються розклади в ряди}~--- автор вважає, що читач розуміє, як вони працюють.
\end{enumerate}

\textbf{Порівняно з Калитою:} там теж є пропуски, але вони іншого типу~--- автор пропускає \emph{фізичні пояснення}, а не математичні викладки. Наприклад, у Калити формула
\[
\vec{E} = \frac{1}{4\pi\varepsilon_0} \frac{q\vec{e}_r}{r^2}
\]
подана без виведення, але з поясненням, що таке \(\vec{e}_r\) і як напрямлений вектор. У Пономаренка формула
\[
\mathbf{E} = \frac{3(\mathbf{p} \cdot \mathbf{r})\mathbf{r}}{r^5} - \frac{\mathbf{p}}{r^3}
\]
подана з виведенням, але проміжні кроки скорочені.

\textbf{Рекомендація:} якщо ви читаєте Пономаренка і натрапляєте на таке місце, спробуйте самостійно відтворити пропущені кроки. Якщо не вдається~--- зверніться до Калити або до класичних підручників (Савельєв, Сивухін), де ці викладки розписані детальніше.

\end{document}